% !TEX root = ../main.tex
\chapter{Linear Algebra}
\label{ch:linalg}

Linear algebra is the mathematical substrate of most machine
learning. Matrix multiplication is the dominant flop; matrix
factorisations underlie PCA, low-rank fine-tuning, gradient
preconditioning, and a variety of regularisation techniques;
linear solvers appear inside differential-equation and
control-theoretic models. \Lucid\ exposes a comprehensive
linear-algebra subpackage at \code{lucid.linalg} that wraps the
LAPACK routines provided by \Accelerate.

This chapter describes the subpackage's surface, the
implementation choices, and the unique constraint that defines
its dispatch behaviour: the \emph{linalg carve-out} introduced in
\chref{ch:backend_dispatch}~Section~\ref{sec:backend_dispatch:carveouts},
which routes every linalg call through the CPU stream regardless
of the user's device.

\section{Overview}
\label{sec:linalg:overview}

The \code{lucid.linalg} surface mirrors \refframework's
\code{torch.linalg} closely, with the same routine names,
signatures, and return conventions. The subpackage contains
roughly 31 functions, falling into five families:

\begin{itemize}
  \item \textbf{Decompositions}: \code{cholesky}, \code{qr},
    \code{svd}, \code{eigh}, \code{lu}, \code{ldl}.
  \item \textbf{Solvers}: \code{solve}, \code{lstsq},
    \code{triangular\_solve}, \code{cholesky\_solve}.
  \item \textbf{Matrix functions}: \code{inv}, \code{det},
    \code{slogdet}, \code{matrix\_exp}, \code{matrix\_power},
    \code{matrix\_rank}.
  \item \textbf{Norms}: \code{norm}, \code{matrix\_norm},
    \code{vector\_norm}, \code{cond}.
  \item \textbf{Products}: \code{cross}, \code{outer},
    \code{vecdot}, \code{multi\_dot}.
\end{itemize}

\subsection{The H8 namespace rule}
\label{ssec:linalg:h8_rule}

Per Hard Rule H8, every linalg op is accessible only via
\code{lucid.linalg.*}, never at the top level
(\code{lucid.cholesky} is invalid) and never as a Tensor method
(\code{x.norm()} is invalid). The rule disambiguates: \code{norm}
could mean vector norm or matrix norm, and the subpackage
namespace makes the intent explicit.

The rule also catches user errors at API surface. A user typing
\code{x.norm()} gets an AttributeError instead of an undocumented
default; the message points to \code{lucid.linalg.norm} or
\code{lucid.linalg.vector\_norm} per the user's actual intent.

\section{The Linalg Carve-Out, Restated}
\label{sec:linalg:carveout}

Linear-algebra decompositions are not implemented in \Metal.
\MLX\ provides them only on its CPU stream, where they delegate
to \Accelerate's LAPACK. The dispatcher recognises this and
routes every linalg call through the CPU stream regardless of the
user's selected device.

\subsection{What the user sees}
\label{ssec:linalg:user_view}

\code{lucid.linalg.cholesky(x)} on a \code{metal} tensor:

\begin{enumerate}
  \item Dispatcher recognises linalg op, applies carve-out.
  \item \code{SharedStorage} bridge transfers \code{x} to CPU
    view (zero-copy in the contiguous case).
  \item \MLX\ runs the decomposition on its CPU stream, calling
    \Accelerate's \code{spotrf}.
  \item Result wrapped back as a metal tensor.
\end{enumerate}

The bridge and the wrap-back are sub-microsecond; the
decomposition cost dominates. For matrices larger than $64 \times
64$, the cost is determined entirely by the LAPACK routine, and
the choice of device is fictive.

\subsection{Why this is acceptable}
\label{ssec:linalg:why_acceptable}

LAPACK on Apple's AMX coprocessor through \Accelerate\ is
competitive with the best open-source implementations and faster
than any naive Metal kernel for decompositions of any non-trivial
size. The carve-out is the right architectural choice given the
hardware reality.

\section{Decomposition Routines}
\label{sec:linalg:decompositions}

\subsection{Cholesky}
\label{ssec:linalg:cholesky}

\code{cholesky(A)} returns the lower-triangular Cholesky factor
$L$ such that
\[
A \;=\; L L^\top, \qquad L_{ii} > 0,
\]
requiring $A$ symmetric positive definite. The factor is unique.
The routine is the fastest decomposition: it costs
$N^3/3 + O(N^2)$ flops, half of LU's cost.

The algorithm computes $L$ column by column:
\[
L_{jj} \;=\; \sqrt{A_{jj} - \sum_{k<j} L_{jk}^2},
\qquad
L_{ij} \;=\; \frac{1}{L_{jj}}\!\left(A_{ij} - \sum_{k<j} L_{ik} L_{jk}\right)
\quad (i > j).
\]
A non-positive value under the square root signals
non-positive-definite $A$; the routine raises.

Backward: given $\bar L$ (the cotangent of $L$), define
\[
P \;=\; \operatorname{tril}(L^\top \bar L) - \tfrac{1}{2}\operatorname{diag}(L^\top \bar L).
\]
Then
\[
\bar A \;=\; \tfrac{1}{2}\, L^{-\top} (P + P^\top) L^{-1}.
\]
\Lucid's autograd registers this analytical gradient; the user
does not see the derivation.

\subsection{QR}
\label{ssec:linalg:qr}

\code{qr(A)} returns $(Q, R)$ where $A = QR$, $Q$ is orthogonal,
$R$ is upper triangular. Two modes:

\begin{itemize}
  \item \code{mode='reduced'} (default): returns
    $Q \in \mathbb{R}^{m \times k}$, $R \in \mathbb{R}^{k \times
    n}$ for $A \in \mathbb{R}^{m \times n}$, $k = \min(m, n)$.
  \item \code{mode='complete'}: returns full $Q \in
    \mathbb{R}^{m \times m}$.
\end{itemize}

Implementation: Householder reflections via \Accelerate's
\code{sgeqrf} for the factorisation, \code{sorgqr} for the
explicit $Q$ reconstruction.

\subsection{SVD}
\label{ssec:linalg:svd}

\code{svd(A)} returns $(U, S, V^\top)$ where $A = U \Sigma V^\top$
and $\Sigma = \text{diag}(S)$. The implementation uses
\Accelerate's \code{sgesdd} (divide-and-conquer), which is
asymptotically faster than the QR-based \code{sgesvd} for
matrices larger than a few hundred entries on each side.

Modes:
\begin{itemize}
  \item \code{full\_matrices=True}: $U \in \mathbb{R}^{m \times
    m}$, $V \in \mathbb{R}^{n \times n}$, $S \in \mathbb{R}^{\min(m,
    n)}$.
  \item \code{full\_matrices=False}: $U \in \mathbb{R}^{m \times
    k}$, $V \in \mathbb{R}^{n \times k}$, $k = \min(m, n)$.
\end{itemize}

The truncated form is typically what's wanted.

\subsection{Eigendecomposition}
\label{ssec:linalg:eigh}

\code{eigh(A)} returns $(w, V)$ where $A v_i = w_i v_i$ for $A$
symmetric (Hermitian). Implementation: \code{ssyevd}
(divide-and-conquer).

\Lucid\ does not currently expose \code{eig} for general
non-symmetric matrices, on the grounds that ML rarely needs it
and \refframework\ deprecated it.

\section{Solvers}
\label{sec:linalg:solvers}

\subsection{Solve}
\label{ssec:linalg:solve}

\code{solve(A, B)} solves $AX = B$ for $X$. Implementation: LU
factorisation via \code{sgesv}.

\subsection{Least squares}
\label{ssec:linalg:lstsq}

\code{lstsq(A, B)} solves $\min_X \|AX - B\|_F$ via SVD-based
divide-and-conquer (\code{sgelsd}). Returns the minimum-norm
solution when $A$ is rank-deficient.

\subsection{Triangular solve}
\label{ssec:linalg:triangular_solve}

\code{triangular\_solve(A, B, upper=True)} solves $AX = B$ when
$A$ is triangular. Cheaper than general solve: $O(N^2)$ flops
instead of $O(N^3)$.

\section{Matrix Functions}
\label{sec:linalg:matrix_functions}

\subsection{inv}
\label{ssec:linalg:inv}

\code{inv(A)} returns $A^{-1}$. Implementation: LU + back
substitution (\code{sgetrf} + \code{sgetri}).

For solving systems, prefer \code{solve(A, B)} over \code{inv(A) @
B}: more numerically stable and faster.

\subsection{Determinant}
\label{ssec:linalg:det}

\code{det(A)} returns $\det(A)$. For ill-conditioned matrices,
prefer \code{slogdet(A)} which returns $(\text{sign},
\log|\det|)$ and avoids overflow / underflow.

\subsection{Matrix exponential}
\label{ssec:linalg:matrix_exp}

\code{matrix\_exp(A)} returns $e^A = \sum_{k} A^k / k!$.
Implementation: Pade approximation with scaling and squaring.

\section{Backward Through Linalg}
\label{sec:linalg:backward}

The backward for each linalg routine has a known closed form
that \Lucid\ has implemented. The derivations are textbook but
not always pleasant; \tabref{tab:linalg:backward_summary}
summarises the form.

\begin{table}[t]
\centering
\small
\begin{tabular}{ll}
\toprule
Routine & Backward form \\
\midrule
cholesky  & triangular solve involving $L^{-1}$ \\
qr        & dense formula with $Q^\top \bar Q$, $R^{-1}$ \\
svd       & three terms; depends on $\bar U$, $\bar S$, $\bar V$ \\
eigh      & two terms; involves $V^\top \bar V$, eigenvalue gaps \\
solve     & two solves: $X^\top$, then propagate to $\bar A$ \\
det       & $\det(A) \cdot (A^{-1})^\top \cdot \bar y$ \\
slogdet   & $(A^{-1})^\top \cdot \bar y_{\log}$ \\
matrix\_exp & series expansion of $e^A \bar Y$ \\
\bottomrule
\end{tabular}
\caption[Linalg backward forms.]{Closed-form backward
expressions for the main linalg routines. The expressions are
generally well-behaved at well-conditioned inputs and become
numerically delicate at near-singular ones; the autograd engine
returns NaN in the latter case.}
\label{tab:linalg:backward_summary}
\end{table}

\subsection{Failure modes}
\label{ssec:linalg:failure_modes}

The backward through SVD or eigh becomes singular when the input
has repeated singular values / eigenvalues. The formula has
$1 / (\sigma_i - \sigma_j)$ terms; equal singular values produce
division by zero.

\Lucid's implementation returns NaN in these cases, matching
\refframework's behaviour. The mitigation is to add a small
random perturbation to the input.

\section{Performance Notes}
\label{sec:linalg:perf}

Measured throughput on M3~Max for representative sizes:

\begin{table}[t]
\centering
\small
\begin{tabular}{lrr}
\toprule
Op & Size & Time \\
\midrule
cholesky & $1024 \times 1024$  & 14~ms \\
cholesky & $4096 \times 4096$  & 410~ms \\
qr       & $1024 \times 1024$  & 28~ms \\
qr       & $4096 \times 4096$  & 920~ms \\
svd      & $1024 \times 1024$  & 110~ms \\
svd      & $4096 \times 4096$  & 4.8~s \\
eigh     & $1024 \times 1024$  & 96~ms \\
eigh     & $4096 \times 4096$  & 4.2~s \\
inv      & $1024 \times 1024$  & 22~ms \\
\bottomrule
\end{tabular}
\caption[Linalg throughput on M3~Max.]{Measured wall-clock times
for square-matrix decompositions on M3~Max via \Accelerate's
LAPACK. Cholesky is the cheapest (positive-definite assumption);
SVD and eigh are the most expensive (dense decomposition with
iterative refinement).}
\label{tab:linalg:perf}
\end{table}

\section{QR and Householder Reflections}
\label{sec:linalg:qr_detail}

The QR factorisation is the workhorse of orthogonal-projection
methods (least squares, Gram--Schmidt orthogonalisation, the QR
algorithm for eigenvalues) \cite{golub_van_loan_2013,trefethen_bau_1997}.
\Lucid\ implements it via \Accelerate's Householder routine
\code{sgeqrf} \cite{householder_1958}, which is the standard
numerically-stable construction.

\subsection{Householder reflections}
\label{ssec:linalg:householder}

A Householder reflection $H = I - 2 v v^\top / (v^\top v)$ for a
unit vector $v$ is a symmetric orthogonal matrix that reflects
any input through the hyperplane orthogonal to $v$. The key
property: given any vector $x$, we can choose $v$ so that $H x$
is a multiple of $e_1$ (the first standard basis vector):
\[
v \;=\; x - \mathrm{sign}(x_1)\,\|x\|\,e_1.
\]

QR proceeds by applying Householder reflections to the columns of
$A$ in sequence, each reflection zeroing out the entries below
the diagonal of one column:
\[
H_n \cdots H_2 H_1 A \;=\; R \quad (\text{upper triangular}),
\]
and $Q = H_1 H_2 \cdots H_n$ is the product of reflections. The
algorithm costs $\sim 2 m n^2 - 2 n^3 / 3$ flops for
$A \in \mathbb{R}^{m \times n}$ with $m \geq n$.

\subsection{The pseudo-inverse via QR}
\label{ssec:linalg:pinv_qr}

The Moore--Penrose pseudo-inverse of $A \in \mathbb{R}^{m \times
n}$ (with $m \geq n$, full column rank) is
\[
A^+ \;=\; (A^\top A)^{-1} A^\top \;=\; R^{-1} Q^\top,
\]
the second form using $A = QR$. The computation via QR is
numerically preferable to forming and inverting $A^\top A$
because $\kappa(A^\top A) = \kappa(A)^2$ (the condition number
squares).

\Lucid's \code{linalg.pinv} uses the QR-based form for tall
matrices and the SVD-based form for general rectangular and
rank-deficient inputs.

\subsection{Backward through QR}
\label{ssec:linalg:qr_backward}

The backward through QR is non-trivial. Given $A = QR$ and
cotangents $\bar Q, \bar R$, the input cotangent is
\[
\bar A \;=\; \bigl[Q\,(\bar R\, R^\top + R\, \bar R^\top) +
(I - Q Q^\top)\, \bar Q\bigr] R^{-\top}
\]
with the symmetric correction
\[
\bar R\, R^\top + R\, \bar R^\top
\;\text{replaced by}\;
\operatorname{copyltu}(\bar R\, R^\top - Q^\top \bar Q),
\]
where \code{copyltu} copies the lower triangle to the upper to
enforce symmetry. The derivation requires careful tracking of the
implicit constraint that $Q^\top Q = I$. \Lucid's autograd ships
the correct backward; the derivation is in the source comments of
\code{lucid/\_C/ops/linalg/qr\_backward.cpp}.

\section{SVD in Depth}
\label{sec:linalg:svd_detail}

The singular value decomposition factors any $A \in
\mathbb{R}^{m \times n}$ as
\[
A \;=\; U \Sigma V^\top, \quad
U \in \mathbb{R}^{m \times m},\ V \in \mathbb{R}^{n \times n},\
\Sigma \in \mathbb{R}_{\geq 0}^{m \times n}\ \text{diagonal}.
\]
The columns of $U$ are the left singular vectors, the columns of
$V$ are the right singular vectors, and the diagonal entries of
$\Sigma$ (ordered decreasingly) are the singular values
$\sigma_1 \geq \sigma_2 \geq \cdots \geq \sigma_{\min(m,n)} \geq
0$.

\subsection{Geometric interpretation}
\label{ssec:linalg:svd_geom}

The SVD says any linear map decomposes into a rotation $V^\top$,
a scaling along axes $\Sigma$, and another rotation $U$. The
singular values measure how much each axis is stretched; the
right singular vectors are the input directions of maximum
stretch; the left singular vectors are the output directions.

\subsection{Low-rank approximation}
\label{ssec:linalg:low_rank}

The Eckart--Young theorem \cite{eckart_young_1936} says the best
rank-$k$ approximation of $A$ in the Frobenius norm is obtained
by truncating the SVD:
\[
A_k \;=\; \sum_{i=1}^{k} \sigma_i u_i v_i^\top,
\]
with error
\[
\|A - A_k\|_F^2 \;=\; \sum_{i > k} \sigma_i^2.
\]
This is the mathematical basis of PCA (with mean-centred data),
of low-rank fine-tuning (LoRA), and of randomised matrix sketching.

\subsection{Backward through SVD}
\label{ssec:linalg:svd_backward}

The SVD backward involves three cotangents ($\bar U$, $\bar
\Sigma$, $\bar V$) producing $\bar A$. The expression has a
delicate term
\[
F_{ij} \;=\; \frac{1}{\sigma_j^2 - \sigma_i^2}, \quad i \neq j,
\]
which becomes singular when any two singular values are equal.
The full backward formula is
\begin{align*}
\bar A \;=\;
&U\,\bar\Sigma\, V^\top \\[2pt]
&+ \tfrac{1}{2} U (F \odot (U^\top \bar U - \bar U^\top U)) \Sigma V^\top \\
&+ \tfrac{1}{2} U \Sigma (F \odot (V^\top \bar V - \bar V^\top V)) V^\top.
\end{align*}
The $F$ terms vanish when only $\bar \Sigma$ is non-zero (e.g.,
when backpropagating only through singular values), which is the
case for nuclear-norm regularisation and is therefore numerically
safe. The $F$ terms become NaN when the input has repeated
singular values; \Lucid\ returns NaN in that case and the user
must regularise.

\section{Eigendecomposition: eigh}
\label{sec:linalg:eigh_detail}

For a symmetric (Hermitian) matrix $A \in \mathbb{R}^{n \times
n}$, the spectral theorem gives
\[
A \;=\; V \Lambda V^\top, \quad
V \in \mathbb{R}^{n \times n}\ \text{orthogonal},\
\Lambda = \operatorname{diag}(\lambda_1, \ldots, \lambda_n).
\]

\Lucid's \code{eigh} returns $(\lambda, V)$ ordered with
$\lambda_1 \leq \cdots \leq \lambda_n$.

\subsection{Backward through eigh}
\label{ssec:linalg:eigh_back}

Given cotangents $\bar \lambda$ and $\bar V$, the input cotangent
is
\[
\bar A \;=\; V\,\bigl(\bar\Lambda + (G \odot (V^\top \bar V))\bigr) V^\top,
\]
where $\bar\Lambda = \operatorname{diag}(\bar\lambda)$ and
\[
G_{ij} \;=\; \begin{cases}
\dfrac{1}{\lambda_j - \lambda_i} & i \neq j \\
0 & i = j.
\end{cases}
\]
Like the SVD backward, the $G$ term diverges at repeated
eigenvalues. The same regularisation advice applies.

\subsection{Inverse iteration for selected eigenvalues}
\label{ssec:linalg:inverse_iteration}

For very large symmetric matrices where only a few extreme
eigenvalues are wanted (the typical case in spectral clustering
or graph Laplacian computations), full \code{eigh} is wasteful.
\Lucid\ does not currently expose a partial-eigendecomposition
routine; users needing one resort to power iteration or
\code{scipy.sparse.linalg.eigsh} for sparse matrices.

\section{Cross, Outer, vecdot, and multi-dot}
\label{sec:linalg:products}

The non-decomposition products in \code{lucid.linalg}.

\subsection{cross}
\label{ssec:linalg:cross}

The cross product
\[
a \times b \;=\; \bigl(a_2 b_3 - a_3 b_2,\ a_3 b_1 - a_1 b_3,\
a_1 b_2 - a_2 b_1\bigr)
\]
for $a, b \in \mathbb{R}^3$, with batch broadcasting. Backward:
$\bar a = \bar y \times b$, $\bar b = a \times \bar y$ (using
the antisymmetric identity).

\subsection{outer}
\label{ssec:linalg:outer}

\code{outer(a, b)} returns $a b^\top \in \mathbb{R}^{m \times
n}$ for $a \in \mathbb{R}^m$, $b \in \mathbb{R}^n$. The same as
\code{a.unsqueeze(-1) * b.unsqueeze(-2)} but with the explicit
linalg semantic.

\subsection{vecdot}
\label{ssec:linalg:vecdot}

\code{vecdot(a, b, dim=-1)} returns the inner product
$\sum_i a_i b_i$ along the named axis. Backward is
$(b\,\bar y,\ a\,\bar y)$.

\subsection{The multi-dot chain product}
\label{ssec:linalg:multi_dot}

\code{multi\_dot([A1, A2, ..., An])} computes the product $A_1
A_2 \cdots A_n$ in the order that minimises total flops. The order
matters substantially: for a chain $M \cdot v$ where $M$ is
$n \times n$ and $v$ is $n \times 1$, computing right-to-left
costs $n^2$ flops, while left-to-right would cost $n^3$ if done
naively.

The minimum-cost parenthesisation is the classic dynamic-
programming problem with $O(n^3)$ table-filling cost in the
number of matrices.

\section{Numerical Conditioning}
\label{sec:linalg:conditioning}

The numerical behaviour of a linear-algebra routine depends on
the \emph{condition number} of its input. For a square invertible
matrix $A$, the 2-norm condition number is
\[
\kappa_2(A) \;=\; \|A\|_2 \cdot \|A^{-1}\|_2 \;=\; \frac{\sigma_{\max}(A)}{\sigma_{\min}(A)}.
\]
A small perturbation $\delta A$ in the input can cause a
perturbation in the output bounded by $\kappa(A) \cdot \|\delta
A\| / \|A\|$. Routines on poorly-conditioned matrices lose
precision proportionally.

\subsection{Recognising ill-conditioning}
\label{ssec:linalg:ill_cond}

\code{cond(A)} returns $\kappa_2(A)$ via the SVD. Heuristics:
\begin{itemize}
  \item $\kappa < 10^3$: well-conditioned; full FP32 precision.
  \item $10^3 < \kappa < 10^6$: moderately ill-conditioned; expect
    3--4 digit precision loss in FP32.
  \item $\kappa > 10^9$: ill-conditioned; consider FP64 or
    regularisation.
\end{itemize}

\subsection{Regularisation}
\label{ssec:linalg:regularisation}

The standard remedy for ill-conditioning is to add a small
multiple of the identity:
\[
A' \;=\; A + \lambda I.
\]
The condition number is bounded:
\[
\kappa_2(A') \;\leq\; \frac{\sigma_{\max}(A) + \lambda}{\sigma_{\min}(A) + \lambda}
\;\leq\; \frac{\sigma_{\max}(A)}{\lambda}.
\]
The trade-off is bias: $A'$ is no longer the original matrix.
Typical regularisation strengths: $\lambda \sim 10^{-6}$ for
mild conditioning issues.

\section{Implementation Mapping}
\label{sec:linalg:impl_mapping}

\tabref{tab:linalg:impl_mapping} maps the public functions to
their LAPACK backends.

\begin{table}[t]
\centering
\small
\begin{tabular}{lll}
\toprule
Lucid function & LAPACK routine(s) & Notes \\
\midrule
\code{cholesky}        & \code{spotrf}, \code{dpotrf}   & SPD only \\
\code{qr}              & \code{sgeqrf}, \code{sorgqr}   & Householder \\
\code{svd}             & \code{sgesdd}                  & divide-and-conquer \\
\code{eigh}            & \code{ssyevd}                  & divide-and-conquer \\
\code{solve}           & \code{sgesv}                   & LU with partial pivoting \\
\code{lstsq}           & \code{sgelsd}                  & SVD-based DC \\
\code{inv}             & \code{sgetrf} + \code{sgetri}  & LU + back-substitution \\
\code{det}             & \code{sgetrf}                  & product of diagonal \\
\code{slogdet}         & \code{sgetrf}                  & log-product of diagonal \\
\code{matrix\_exp}     & custom Pade + scaling          & no direct LAPACK \\
\code{triangular\_solve} & \code{strtrs}                & specialised solve \\
\code{cholesky\_solve} & \code{spotrs}                  & specialised solve \\
\code{ldl}             & \code{ssytrf}                  & symmetric indefinite \\
\bottomrule
\end{tabular}
\caption[Linalg backend mapping.]{Mapping from \Lucid's
\code{linalg} functions to the underlying LAPACK routines that
\Accelerate\ provides. The CLAPACK \code{LAPACKE\_*} wrappers are
used throughout for the row-major-aware API.}
\label{tab:linalg:impl_mapping}
\end{table}

\section{Applications in ML}
\label{sec:linalg:applications}

A few patterns where the linalg subpackage appears in modern ML.

\subsection{Principal Component Analysis (PCA)}
\label{ssec:linalg:pca}

Given data $X \in \mathbb{R}^{n \times d}$ (mean-centred), PCA
finds the directions of maximum variance:
\[
X = U \Sigma V^\top,
\quad \text{principal components} = V_{:,1:k},
\quad \text{scores} = U_{:,1:k} \Sigma_{1:k,1:k}.
\]
In \Lucid:
\begin{lstlisting}[style=lucid,language=LucidPython]
X = X - X.mean(dim=0, keepdims=True)
U, S, Vt = lucid.linalg.svd(X, full_matrices=False)
components = Vt[:k]
scores = U[:, :k] * S[:k]
\end{lstlisting}

\subsection{Low-rank fine-tuning (LoRA)}
\label{ssec:linalg:lora}

LoRA \cite{hu_lora_2021} represents a weight delta as a low-rank
product:
\[
\Delta W = B A, \quad B \in \mathbb{R}^{d \times r},\
A \in \mathbb{R}^{r \times d},\ r \ll d.
\]
The fine-tuned weight is $W + B A$. The original weight $W$ is
frozen; only $A$ and $B$ are trained. \Lucid\ provides this as
a module; the underlying linear algebra is plain \code{matmul}.

\subsection{Spectral normalisation}
\label{ssec:linalg:spec_norm}

The spectral norm of a weight matrix $\|W\|_2 = \sigma_1(W)$
bounds the Lipschitz constant of the linear map. Spectral
normalisation \cite{miyato_specnorm_2018} divides $W$ by an
estimate of its largest singular value, computed via power
iteration (one or two iterations per forward) rather than a full
SVD.

\section{Summary and Forward References}
\label{sec:linalg:summary}

\code{lucid.linalg} exposes a comprehensive set of linear-algebra
routines that delegate to \Accelerate's LAPACK via the CPU
stream. The carve-out makes device choice fictive; the bridge
cost is negligible. Backward forms are known in closed form for
every routine, with the caveat that SVD and eigh become
numerically singular at degenerate inputs.

The next chapter, \chref{ch:fft}, treats the FFT subpackage that
shares the same dispatch carve-out.

\begin{lucidnote}
For the user: \code{lucid.linalg} is the only namespace. Device
argument is informational only for linalg ops; the actual
computation always runs on CPU through \Accelerate.
\end{lucidnote}
